\documentclass[a4paper]{seasonalletter}
\SeasonalSetup{season=autumn,weekday=wednesday,paper=grid,
  label={A letter from the reading room}}
\begin{document}
\begin{SeasonalLetter}
\SeasonalHeading{The quiet art of keeping in touch}
\hfill 5 November 2026\par
\SeasonalOpening{Dear friend,}
A letter gives a small part of the day somewhere to stay. This page is a working
LaTeX document: its words remain selectable, its paragraphs break naturally,
and its decoration is drawn as vector artwork behind the text.

The grid is deliberately quiet. Its lines belong to the page, so they remain
in place as the letter grows. The oak leaves stay in the margin, leaving the
centre clear for reading.

This afternoon the light reached the edge of the desk and stopped beside an
unfinished book. I thought you would have liked the silence of that moment.
There was nothing urgent to report, which seemed a good reason to write.
\SeasonalClosing{With warm wishes,}{A fellow reader}
\newpage
\SeasonalHeading{A little more room}
This is the continuation of the same letter. The package keeps the border and
a small weekday emblem, while leaving out the large corner motifs and label.
The page number belongs to this letter and does not depend on the document's
printed page counter.

For an undecorated continuation, choose \texttt{continuation=none}. For the
same ornament on every page, choose \texttt{continuation=full}. The paper
colour and optional ruling continue in either case.

The end of a letter should have as much room as its beginning. A quiet second
page makes that possible without requiring the writer to position paragraphs
by hand.
\SeasonalClosing{Until next time,}{A fellow reader}
\end{SeasonalLetter}
\end{document}
